\section{Verification and reproducibility}\label{sec:verification}
The final release checks are organized by mathematical obligation.
\begin{itemize}
\item The compact checker authenticates $105$ inherited source files and
all four fixed numerical payloads. For each of the five circuits it checks
topology, liveness, the complete parent/child gate relation, and every
polynomial stage through $32$--$37$ exact cuts. Nine auxiliary identities
include the strong-trace algebra and polynomial completeness map.
\item It computes the complete highest homogeneous forms. Eighteen
deliberate operator corruptions per circuit are rejected, for a total of
$90$. Normal and optimized-Python receipts agree. Exact reconstruction
of the fixed tuple and byte-for-byte circuit rebuilding pass.
\item The direct checker independently represents monomials as
variable/exponent pairs. It checks the actual diagonal-power source edit,
$18$ Horn identities, $202$ predecessor monic projections, the selected
affine maps, and all six new radical projections for each endpoint.
\item It independently expands every final quartic, verifies every
coefficient, checks all reported metrics, and verifies the terminal
Hessian ranks modulo $101$. Five deliberate corruptions of directions,
restorations, coordinate deletion, residuals, and expanded coefficients
are rejected.
\item The scoped predecessor checks pass both normally and with Python
assertions disabled. A disposable reconstruction reproduces $28$
selected mathematical JSON files byte for byte. These retained
dependency subsets omit historical PDFs and obsolete experiments;
their scoped entry points are explicitly documented.
\item A root release replay checks the component verifiers and ensures
that the report's literal residual appendix, comparison tables, and
degree examples match their canonical data. The PDF is rendered and
visually inspected before packaging.
\end{itemize}

The release uses Python 3.10 or newer and its standard library for
these exact checks. Run from the extracted package directory:
\begin{verbatim}
python3 -B verify.py
python3 -B verify.py --rebuild
sh build_report.sh
\end{verbatim}
The first command checks the frozen artifacts without changing them.
The optional second additionally regenerates the selected artifacts
in disposable locations. The last rebuilds the paper from its TeX
source. Internal independent mathematical reviews are retained in
\texttt{review/}; they are AI-assisted reviews, not external human
peer review.


The canonical equations use ordered sparse terms:
$[a,[i_1,\ldots,i_r]]$ denotes $a x_{i_1}\cdots x_{i_r}$.
Repeated indices are repeated factors, and an empty index list is a
constant. Every array is finite. The evaluator accepts an ordered tuple
of exact integers; it is not a solver or a consistency decision procedure.

Every compact coefficient is included with its inherited exact
reconstruction. The actual native tuple, not an arbitrary tuple
satisfying a few bounds, is bound to the consistency statement.
Checks distinguish coefficient reconstruction, literal circuit identity,
degree certification, expansion, transformation correctness, and the
unbounded mathematical proof.

The source snapshot and archive hashes are in the ledger. The final
manifest authenticates the deliverable bytes. The TeX source builds
with a standard TeX Live installation; font versions and PDF metadata
can change the PDF hash without altering the canonical equations.

\section{What this iteration adds, and further research}

The compact advance is the use of the explicit strong-trace congruence
$f\equiv1\pmod c$, together with a complete signed proof, to replace
$f^2(cT-R)-c$ by $f(cT-R)-c$. This lowers degree at the same binary operation count: the new trace
is evaluated as $f=a+(1+a)$ using the already shared $1+a$. The new polynomial parent-to-child map gives a further
algebraic check on completeness.

The direct advance is a reusable integral projection mechanism that
finds six hidden removable directions. The $t^t$ lemma separately
simplifies an actual power request and its justification. The projection
applies to more than one inherited source, and the label lemma is
proved before any quadratic projection.

The following questions are concrete directions for subsequent work.
No unproved answer is used in this release.
\begin{enumerate}
\item Can a weaker auxiliary congruence yield a further compact
trace deletion? A new index argument would be needed: the present
residue modulo $f$ does essential work.
\item Can a changed residual presentation expose more removable
directions after the common radical becomes zero? Constant invertible
row operations preserve the common radical, so they cannot suffice alone.
\item Can a structured nonlinear coordinate change preserve quadratic
residuals and expose more unit coordinates? A valid proposal needs an
integral inverse or an explicit existential restoration.
\item Can a simpler proof calculus reduce the direct front end enough
to compensate for its arithmetic overhead and calibration? Counting
rule schemata alone does not measure the final polynomial.
\item Can the synchronized power compiler use the diagonal nature
of $F=t^t$ more deeply than changing its base? Soundness must still
cover arbitrary integer assignments.
\item Can the radical-projection and trace-congruence lemmas first
be formalized as reusable components, then connected to the internal-ZFC
calibration? This would address the actual formal gap rather than
recheck only finite numerical cases.
\end{enumerate}
